\section{Theorem C: regime arithmetic, boundary obstructions, and a second certificate family}\label{sec:C}
\textbf{Full proposed parameter-hierarchy theorem status: Open (not proved).}

We investigate the empirical starting-position template
\begin{equation}\label{eq:template}
\mathcal{T}(k,A):\qquad P_0(n)\iff n\bmod p\in R,\qquad
p=2k+A+2,\quad R=\{0,A,k+A+1\},
\end{equation}
where $k,A$ are positive integers and $P_0(n)$ means that the unrestricted start with $n$ counters is losing. These observations concern this game's specific recurrence, not a claim previously established in the cited literature. The background on subtraction games and no-repetition variants motivates the question \cite{winningways2001,gurvich2012}.

\subsection{Uniform analytic results: proved}
\begin{proposition}[Immediate upper bound on any valid template]\label{prop:Cupper}
If $\mathcal{T}(k,A)$ holds for $D(k,1)$, then necessarily $k<A$.
\end{proposition}
\begin{proof}
Equation~\eqref{eq:template} declares $n=A$ a losing unrestricted start, since $A\in R$. If $k\ge A$, however, removing exactly $A$ counters is a legal \emph{first} move from $n=A$ and leaves an empty heap, where the next player has no legal move. Therefore the first player wins from $n=A$, a contradiction.
\end{proof}

\begin{proposition}[Arithmetic collapse of the three-residue template]\label{prop:Ccollapse}
Let $p=2k+A+2$ and $R=\{0,A,k+A+1\}\pmod p$ with $k,A\ge1$. The binary $p$-periodic indicator of $R$ has a smaller positive period if and only if $k=A-1$. In that case $p=3A$, $R=\{0,A,2A\}$, and the least positive period is $A$.
\end{proposition}
\begin{proof}
The residues satisfy $0<A<k+A+1<p$, so there are exactly three. If the indicator has a smaller period, a nonzero translation of the cyclic group $\mathbb Z/p\mathbb Z$ preserves $R$. Every orbit of a fixed translation has equal size, which must divide $|R|=3$. Since a nonzero translation has no fixed point, its orbit size must equal three. The invariant set is therefore the orbit $\{0,p/3,2p/3\}$. The smaller positive residue is $A$, so $A=p/3$, equivalently $3A=2k+A+2$, or $k=A-1$. Conversely, if $k=A-1$, then $p=3A$ and $R=\{0,A,2A\}$; membership depends only on divisibility by $A$, and $A$ is clearly least because only residue zero is selected modulo $A$.
\end{proof}

For $A=54$ the upper bound allows only $k\le53$, and the endpoint $k=53$ forces an arithmetic period collapse. Both statements are unconditional \emph{implications}; they do not by themselves establish that the template is realized by a game for any particular $k$.

\subsection{Certified second affine family: proved with finite computation}
\begin{theorem}[Second fixed-parameter affine regime]\label{thm:Cfinite}
For every fixed integer $42\le k\le53$, the unrestricted starting position in $D(k,1)$ is losing exactly when
\[\boxed{n\bmod(2k+56)\in\{0,54,k+55\}.}\]
For $42\le k\le52$ the least positive starting period is $2k+56$. For $k=53$ the least positive starting period is $54$.
\end{theorem}
\begin{proof}[Exact finite certificates and induction]
For each of the twelve integers, the independent \texttt{methodology/theorem\_C/verify\_theorem\_c.py} invokes the exact Theorem-A interval recurrence, seeks a repeated window of $k$ P-state rows with displacement $p=2k+56$, and checks the formula through $t+k+p-1$. The $t$ and bridge endpoints are printed in Table~\ref{tab:Ccert}. Theorem-B Lemmas~\ref{lem:Bwindow} and~\ref{lem:Bbridge} (reproved by identical elementary induction if this theorem is read independently) extend those finite observations to every $n$. Each row and starting-position outcome through the bridge is checked against a separate direct Sprague--Grundy solver. Thus the computation supplies explicit verifiable finite premises, while mathematical induction supplies the infinite consequence. Proposition~\ref{prop:Ccollapse} establishes the claimed least positive periods; the only collapse in the listed interval occurs at $k=53=A-1$.
\end{proof}
\subsection{Exact fixed-parameter certificate table}
\begin{longtable}{rrrrrl}
\caption{Theorem C fixed-parameter certificates for $A=54$.}\label{tab:Ccert}\\
\toprule $k$ & candidate $p$ & least & $t$ & bridge end & SG \\
\midrule\endfirsthead\toprule $k$ & candidate $p$ & least & $t$ & bridge end & SG \\
\midrule\endhead
42 & 140 & 140 & 111 & 292 & PASS \\
43 & 142 & 142 & 117 & 301 & PASS \\
44 & 144 & 144 & 117 & 304 & PASS \\
45 & 146 & 146 & 115 & 305 & PASS \\
46 & 148 & 148 & 116 & 309 & PASS \\
47 & 150 & 150 & 112 & 308 & PASS \\
48 & 152 & 152 & 114 & 313 & PASS \\
49 & 154 & 154 & 0 & 202 & PASS \\
50 & 156 & 156 & 0 & 205 & PASS \\
51 & 158 & 158 & 0 & 208 & PASS \\
52 & 160 & 160 & 0 & 211 & PASS \\
53 & 162 & 54 & 0 & 214 & PASS \\
\bottomrule\end{longtable}
\subsection{Nearby-parameter counterexamples: proved finite failures}
A template that agrees with early rows need not repeat indefinitely. The exact recurrence and an independent direct P/N solver give the following first disagreement with the indicated $A$-template. Here ``predicted'' and ``actual'' refer to whether the unrestricted start is a P-position.
\begin{longtable}{rrrrrl}
\caption{Counterexamples to extending the fixed-$A$ formulas.}\label{tab:Cfail}\\
\toprule $k$ & $A$ & first $n$ & predicted P? & actual P? & direct P/N \\
\midrule\endfirsthead\toprule $k$ & $A$ & first $n$ & predicted P? & actual P? & direct P/N \\
\midrule\endhead
41 & 54 & 42 & no & yes & PASS \\
54 & 54 & 54 & yes & no & PASS \\
55 & 54 & 54 & yes & no & PASS \\
56 & 54 & 54 & yes & no & PASS \\
57 & 54 & 54 & yes & no & PASS \\
57 & 86 & 58 & no & yes & PASS \\
79 & 86 & 492 & yes & no & PASS \\
80 & 86 & 167 & yes & no & PASS \\
\bottomrule\end{longtable}
\paragraph{Interpretation.} The $k=54,\ldots,57$ failures of $A=54$ at $n=54$ follow already from Proposition~\ref{prop:Cupper}. The $k=79$, $A=86$ failure at $n=492$ is particularly instructive: several earlier blocks of starting outcomes happen to agree with the proposed formula, but its recurrence ultimately disagrees. One must not infer a theorem from a long numerical scan or a finite streak of residue matches.

\subsection{What remains Open}
\begin{conjecture}[Parameterized regime mechanism---Open]\label{conj:Chierarchy}
There may exist a structured set of constants $A$ and corresponding $k$-intervals for which $\mathcal{T}(k,A)$ is valid, with $A=54$ and $A=86$ representing two examples. An eventual symbolic classification may explain those constants, the actual parameter endpoints, and the complete P-state row intervals.
\end{conjecture}
\statusopen\quad No assertion that an infinite hierarchy exists is proved. The finite certificates do not yield a parameter-uniform inductive formula for all rows. In particular we \emph{have not} derived the numbers $54$ and $86$ from first principles, shown why the $A=86$ regime terminates at $k=78$, or ruled out isolated additional ranges of $k$ with some other $A$.

\subsection{Next verifiable proof obligations}
The next step is to derive a finite, symbolic, piecewise affine formula for $Q_n$ uniformly in $k$ on a regime, prove each branch stable under the exact set recurrence, and justify the allowed parameter inequalities. A certificate-assisted theorem for finitely many $k$ values must be distinguished from such a uniform proof.\par
